\documentclass[a4paper,11pt]{article}
% Gemeinsame Präambel -- Format aus der bereitgestellten Vorlage übernommen.
\usepackage[utf8]{inputenc}
\usepackage[T1]{fontenc}
\usepackage{lmodern}
\usepackage[ngerman]{babel}
\usepackage[margin=2.3cm]{geometry}
\usepackage{enumitem}
\usepackage{booktabs}
\usepackage{tabularx}
\usepackage{xcolor}
\usepackage{fancyhdr}
\usepackage{amsmath}
\usepackage{array}
\usepackage[hidelinks]{hyperref}

\definecolor{zwetschge}{HTML}{5B2A86}
\definecolor{offen}{HTML}{C0392B}

\pagestyle{fancy}
\fancyhf{}
\lhead{Zwätschgenabend -- Pubquiz}
\rhead{HSLU Elektrotechnik}
\cfoot{\thepage}
\setlength{\headheight}{14pt}

\setlist[itemize]{itemsep=0.2em, topsep=0.3em}


\newif\ifloesung
\loesungtrue

\newcommand{\loesung}[1]{%
  \ifloesung\par\smallskip{\color{zwetschge}\textit{Lösung: #1}}\par\medskip\fi}

\long\def\mc#1#2?#3#4#5#6#7{%
  \subsection*{#1}
  #2?
  \begin{enumerate}[label=\Alph*)]
    \item #3
    \item #4
    \item #5
    \item #6
  \end{enumerate}
  \loesung{#7}}

\newcommand{\offenfrage}[3]{%
  \subsection*{#1}
  #2
  \loesung{#3}}

\newcommand{\runde}[1]{\newpage\section*{#1}\addcontentsline{toc}{section}{#1}}

\begin{document}

\begin{center}
  {\LARGE\bfseries Zwätschgenabend -- Fragenkatalog}\\[0.4em]
  {\large HSLU Elektrotechnik}\\[0.8em]
  \textit{Zuerst wissen. Später nur noch verdächtig gut raten.}
\end{center}

\noindent
\textbf{Spielidee:} Die ersten beiden Runden sind noch fachlich vertretbar.
Danach wird es zunehmend wichtig, wer wirklich zuhört, wenn im Labor niemand
mehr zuhört. Pro Frage ein Punkt, bei offenen Fragen entscheidet die
Quizleitung über technisch vertretbare Kreativität.

\runde{1. Noch ganz normale Fachfragen}

\mc{1.1 -- Ein ehrlicher Bildschirm}
Auf dem Oszilloskop füllt genau eine Periode eines Signals vier horizontale
Kästchen. Die Zeitbasis steht auf $5\,\mathrm{ms/div}$. Welche Frequenz hat das
Signal?
{$25\,\mathrm{Hz}$}
{$50\,\mathrm{Hz}$}
{$100\,\mathrm{Hz}$}
{So präzise misst man nur, wenn jemand zuschaut.}
{B -- Eine Periode dauert $20\,\mathrm{ms}$, also $50\,\mathrm{Hz}$.}

\mc{1.2 -- Der Widerstand als Lebensberatung}
Warum bekommt eine LED in einer üblichen Schaltung einen Vorwiderstand?
{Damit sie beim Einschalten nicht sofort sehr persönlich mit dem Netzteil wird.}
{Damit ihre Farbe mathematisch stabil bleibt.}
{Damit der Strom begrenzt wird.}
{Damit man auf dem Schema noch etwas zum Beschriften hat.}
{C -- und A beschreibt die Konsequenz davon erstaunlich gut.}

\mc{1.3 -- Leistung ohne Meinung}
An $12\,\mathrm{V}$ fliesst ein Strom von $2\,\mathrm{A}$. Welche Leistung
wird umgesetzt?
{$6\,\mathrm{W}$}
{$10\,\mathrm{W}$}
{$14\,\mathrm{W}$}
{$24\,\mathrm{W}$}
{D -- $P=U\cdot I$. Das Netzteil darf danach trotzdem kurz warm werden.}

\mc{1.4 -- Die höfliche Sicherung}
Welche Komponente unterbricht den Stromkreis im Fehlerfall am ehesten,
ohne dabei so zu tun, als sei nichts passiert?
{Eine Sicherung}
{Ein Pull-up-Widerstand}
{Ein Quarz}
{Ein Student mit einem sehr langen Schraubenzieher}
{A -- der Schraubenzieher ist keine zugelassene Schutzmassnahme.}

\mc{1.5 -- Digital, aber nicht magisch}
Was bedeuten $3{,}3\,\mathrm{V}$ an einem üblichen digitalen Eingang
\emph{typischerweise}?
{Logisch 1}
{Logisch 0}
{Der Eingang will zuerst ein Datenblatt lesen.}
{Dass die Schaltung sicher mit $5\,\mathrm{V}$ betrieben werden darf.}
{A -- mit dem üblichen Hinweis auf Datenblatt, Versorgung und Logikfamilie.}

\mc{1.6 -- Der Puffer, den alle kennen}
Was macht ein Kondensator direkt bei der Versorgung eines ICs besonders
gern?
{Er puffert schnelle Spannungseinbrüche und Störungen.}
{Er übersetzt C nach Python.}
{Er gibt dem IC emotionale Stabilität.}
{Er erhöht die Frequenz, weil er klein ist.}
{A -- deshalb liegt er auch nicht dekorativ neben dem IC.}

\mc{1.7 -- Der erste Fehler}
Eine Schaltung funktioniert nicht. Welche erste Handlung ist am ehesten
Ingenieurwesen und am wenigsten Wunschdenken?
{Den Code komplett neu schreiben.}
{Versorgung, Polarität und offensichtliche Verbindungen messen.}
{Die Schaltung dreimal aus- und wieder einschalten.}
{Erklären, dass sie zuhause noch funktioniert hat.}
{B -- die anderen Varianten sind danach immer noch verfügbar.}

\offenfrage{1.8 -- Die einzige ehrliche Antwort}
Nenne eine Messung, die du vor dem Satz \textit{«Das Layout ist schuld»}
machen würdest.
{Zum Beispiel Versorgungsspannung, Stromaufnahme, Massebezug, Signal direkt
an der Quelle oder Lötstellen. Jede sinnvolle Messung zählt.}

\runde{2. Signale mit Nebenwirkungen}

\mc{2.1 -- Faltung, endlich praktisch}
Was beschreibt die Faltung eines Eingangssignals mit der Impulsantwort
eines linearen zeitinvarianten Systems?
{Wie das System auf das Eingangssignal reagiert.}
{Wie zwei Signale sich im Labor gegenseitig ignorieren.}
{Eine Methode, Kabel platzsparend zu verlegen.}
{Den Moment, in dem aus Mathematik Gruppenarbeit wird.}
{A -- leider exakt die mathematische Version von «das System macht etwas damit».}

\mc{2.2 -- Der Frequenzgang als Charaktertest}
Ein Tiefpass dämpft hohe Frequenzen. Was passiert mit einem Rechtecksignal,
wenn man ihn ernst nimmt?
{Die Flanken werden runder.}
{Es wird automatisch ein Sinus mit besserem Selbstwertgefühl.}
{Die Frequenz verschwindet aus Respekt.}
{Das Signal beantragt einen Wechsel zu Digital Engineering.}
{A -- genügend starke Filterung kann es allerdings sehr sinusähnlich machen.}

\mc{2.3 -- Der Lüfter lügt}
Ein Lüfter wird viel zu langsam abgetastet und scheint auf dem Bildschirm
rückwärts zu laufen. Wie heisst der Effekt?
{Doppler-Effekt}
{Aliasing}
{Faltung}
{Der Lüfter hat eine Meinung zum Praktikum.}
{B -- das Signal tut nur so, als hätte es eine tiefere Frequenz.}

\mc{2.4 -- Nyquist ohne Ausrede}
Welche minimale Abtastrate braucht ein ideal bandbegrenztes Signal mit einer
höchsten Frequenz von $8\,\mathrm{kHz}$ mindestens?
{$4\,\mathrm{kHz}$}
{$8\,\mathrm{kHz}$}
{$16\,\mathrm{kHz}$}
{Eine Rate, die sich dynamisch genug anfühlt.}
{C -- mindestens das Doppelte der höchsten vorkommenden Frequenz.}

\mc{2.5 -- Bode im Alltag}
Was liest man in einem Bode-Diagramm direkt ab?
{Amplituden- und Phasenverlauf über der Frequenz.}
{Die Wahrscheinlichkeit, dass das Labor noch offen ist.}
{Die Anzahl korrekt angeschlossener Masseklemmen.}
{Die Stimmung einer Rückkopplung.}
{A -- die Stimmung lässt sich höchstens daraus ableiten.}

\mc{2.6 -- Die falsche Sicherheit}
Ein Frequenzgang sieht in der Simulation wunderschön aus, die Hardware
schwingt aber. Was ist die beste erste Vermutung?
{Die Realität hat einen Rechenfehler.}
{Parasitische Effekte, Bauteiltoleranzen oder Aufbau wurden idealisiert.}
{Das Oszilloskop will provozieren.}
{Die Frequenzantwort wurde beleidigt.}
{B -- echte Bauteile und echte Leiterplatten haben leider Nebenrollen.}

\mc{2.7 -- Filterfrage mit Musik}
Welche Schaltung lässt langsame Änderungen passieren und dämpft schnelle?
{Tiefpass}
{Hochpass}
{Bandpass}
{Der Filter im Kaffeeautomaten}
{A -- ein Tiefpass ist praktisch der Türsteher für schnelle Signale.}

\offenfrage{2.8 -- Faltung in einem Satz}
Erkläre Faltung so, dass auch ein Mensch mit drei Stunden Schlaf noch
glaubt, sie verstanden zu haben.
{Sinngemäss: Man betrachtet, wie sich die Impulsantwort des Systems für jeden
Zeitpunkt mit dem Eingang überlagert. Gute, verständliche Bilder zählen.}

\runde{3. Amtliche Erkenntnisse}

\mc{3.1 -- Der Knopf}
Im ET-Labor gilt die Taste \texttt{Autoscale} als sicherheitskritisch. Was passiert laut Gerüchten, wenn sie trotzdem gedrückt wird?
{Das Oszilloskop beantragt eine Weiterbildung in Digital Engineering.}
{Der Frequenzgang wird emotional und beginnt, den Raum zu bewerten.}
{Rolf Mettler kehrt aus dem Ruhestand zurück und stellt den Knopf auf manuell.}
{Die Zentralbahn fährt wegen EMV-Bedenken nur noch im Rechtecksignal.}
{C}
§
\offenfrage{3.2 -- Personalie}
Wer ist «the one and only»?
{Rolf Mettler.}

\mc{3.3 -- Protokollauszug}
In einer Besprechung fällt der Satz: «Man könnte das mitunter direkt
beantworten, oder zuerst überlegen, ob die Antwort nicht noch eine
andere Frage braucht. Oder» Wem wird die Aussage zugeordnet?
{Kilian Schuster}
{Thomas Hunziker}
{Christian Jost}
{Rolf Mettler}
{A}

\mc{3.4 -- Sprachstand}
Die Fachstelle für Mehrsprachigkeit untersucht eine Erklärung, die mit einem
Satz über Signalverarbeitung beginnt und als kurze, funktionierende Lösung
endet. Welche Sprache wurde darin am wahrscheinlichsten gesprochen?
{Python}
{Französisch}
{Latex}
{Die Sprache des Datenblatts}
{A}

\mc{3.10 -- Die Augenbrauenmethode}
Welche Reaktion bedeutet in einem analogen Labor mit hoher Wahrscheinlichkeit,
dass eine Erklärung noch einmal vollständig geprüft werden sollte?
{Eine sehr kleine Veränderung im Gesichtsausdruck der betreuenden Person.}
{Eine neue Version der Aufgabenstellung.}
{Ein besonders freundliches «passt schon».}
{Ein Screenshot des Messwerts.}
{A}

\runde{4. Fakultätsübergreifende Messungen}

\mc{4.1 -- Schwierigkeitsindikator}
Die ET-interne Kommission sucht einen verlässlichen Hinweis darauf, dass ein
Modul nicht sofort zur Flucht zwingt. Welcher Befund genügt?
{Auch WI-Studierende sitzen noch im Raum.}
{Die Aufgabenstellung ist auf Recyclingpapier gedruckt.}
{Der Beamer zeigt ein Balkendiagramm.}
{Jemand hat das Wort «Stakeholder» verwendet.}
{A}

\mc{4.2 -- Etymologie}
Im inoffiziellen Wörterbuch des wahren ETlers wird der Begriff
«Kanonenfutter» nicht erklärt, sondern mit einem einzelnen, fachlich
vollkommen ausreichenden Beispiel illustriert. Welches?
{Gebäudetechniker in ET2.}
{Die Beilage zum Mensamenü.}
{Ein Kollateralschaden beim Praktikumsaufbau.}
{Der Inhalt einer Schublade mit alten Bananensteckern.}
{A}

\end{document}

% Archiv des verworfenen Insider-Entwurfs
\iffalse
\runde{3. Laborfolklore}

\mc{3.1 -- Der Knopf}
Im ET-Labor gilt die Taste \texttt{Autoscale} als sicherheitskritisch. Was passiert laut Gerüchten, wenn sie trotzdem gedrückt wird?
{Das Oszilloskop beantragt eine Weiterbildung in Digital Engineering.}
{Der Frequenzgang wird emotional und beginnt, den Raum zu bewerten.}
{Rolf Mettler kehrt aus dem Ruhestand zurück und stellt den Knopf auf manuell.}
{Die Zentralbahn fährt wegen EMV-Bedenken nur noch im Rechtecksignal.}
{C}

\offenfrage{3.2 -- Die Frage, die keine Erklärung braucht}
Wer ist «the one and only»?
{Rolf Mettler.}

\mc{3.3 -- Der Satz, der länger dauert als die Frage}
Eine Aussage wird mitunter, oder auch nicht, zunächst erklärt; oder man
erklärt mitunter zuerst, weshalb sie vielleicht, oder eben gerade nicht,
erklärt werden muss. Wer hat den Raum betreten?
{Kilian Schuster}
{Thomas Hunziker}
{Giselher Wichmann}
{Ein Chatbot, der zu viel Kaffee hatte}
{A -- wer bei «mitunter» und «oder» bereits die Stimme gehört hat, ist klar im Vorteil.}

\mc{3.4 -- Mehrsprachiges Debugging}
Ein Laptop läuft auf Japanisch, ein Datenblatt wird in einer weiteren Sprache
gelesen, und danach wird ein Signalproblem erstaunlich elegant gelöst. Was
ist in diesem Moment die eigentliche gemeinsame Muttersprache?
{Python}
{Latex}
{C, aber nur wenn ein Semikolon herumliegt}
{Die Sprache, in der Datenblätter nie geschrieben sind}
{A -- bei Thomas Hunziker ist Python keine Programmiersprache, sondern offenbar Grundversorgung.}

\mc{3.5 -- Die Oberarmmessung}
Ein Dozent erscheint mit einer Bizepsstabilität, bei der selbst ein
abgeglichenes Netzteil kurz nachfragt. Welche Untersuchung fordert die
völlig unseriöse Campusforschung?
{Eine Dopingkontrolle}
{Eine EMV-Messung}
{Eine Augenprüfung}
{Ein Update des Stundenplans}
{A -- der betroffene Mikrocontroller-Dozent weiss, wer gemeint ist.}

\mc{3.6 -- Die Augenbraue als Messgerät}
Im Labor hebt sich eine Augenbraue minimal. Welche Handlung hat jetzt die
höchste Erfolgswahrscheinlichkeit?
{Die Messung nochmals von Hand prüfen.}
{Autoscale drücken, damit es endlich objektiv wird.}
{Erklären, dass die Simulation anders aussah.}
{Das Netzteil vorsorglich auf Maximum drehen.}
{A -- ein sehr präziser analoger Kommunikationskanal.}

\mc{3.7 -- Die Kreidetafel als Hochfrequenzgerät}
Eine Erklärung enthält Laplace, Frequenzbereich und Schweizerdeutsch, bis
das Thema plötzlich doch nachvollziehbar wird. Wer ist gemeint?
{Giselher Wichmann}
{Adrian Omlin}
{Erich Styger}
{Patric Eberle}
{A}

\mc{3.8 -- Heavy Metal, aber mit Netzrückwirkung}
Jemand kann einen Raum mit $50\,\mathrm{Hz}$, Leistungselektronik und
«Heavy Metal» füllen, ohne eine einzige Gitarre auszupacken. Wer?
{Adrian Omlin}
{Markus Thalmann}
{Kilian Schuster}
{Thomas Hunziker}
{A}

\mc{3.9 -- Der Bug merkt es selbst}
Ein Fehler im Embedded-Code fühlt sich bereits ertappt, noch bevor jemand
den Debugger geöffnet hat. Wer könnte gerade in der Nähe sein?
{Erich Styger}
{Giselher Wichmann}
{Thomas Graf}
{Ein Pull-down-Widerstand}
{A -- manche Menschen haben eine sehr kurze Distanz zum Fehler.}

\mc{3.10 -- Die harmlose Rückfrage}
Welcher Satz bedeutet mit hoher Wahrscheinlichkeit, dass nun alle ihre
Unterlagen innerlich neu sortieren müssen?
{«Das hättet ihr doch schon gelernt, oder?»}
{«Wir machen heute früher Schluss.»}
{«Das wird nicht geprüft.»}
{«Die Messung sieht plausibel aus.»}
{A -- besonders gefährlich, wenn Patric Eberle dabei freundlich wirkt.}

\mc{3.11 -- Die Verschlüsselung des Alltags}
Ein Hardwareprojekt hat plötzlich einen Quantencomputer, die NSA und ein
ASIC im Raum, obwohl ursprünglich nur ein Bauteil erklärt werden sollte.
Wer ist am wahrscheinlichsten schuld?
{Markus Thalmann}
{Adrian Omlin}
{Christian Jost}
{Der Hausdienst}
{A}

\offenfrage{3.12 -- Nicht im Datenblatt}
Erfinde eine technische Warnmeldung, die im ET-Labor dringend fehlen würde.
{Juryentscheid. Bonus, wenn sie glaubwürdig klingt und trotzdem niemand sie
jemals implementieren dürfte.}

\runde{4. Campuswissen, das man nirgends lernt}

\mc{4.1 -- Kanonenfutter}
Was ist Kanonenfutter?
{Gebäudetechniker in ET2}
{Ein falsch dimensionierter Shunt}
{Der Inhalt einer Kaffeepause um 07:55 Uhr}
{Die erste Zeile eines ungetesteten C-Programms}
{A -- eine Campusweisheit, die ohne weitere Herleitung auskommt.}

\mc{4.2 -- Ein beruhigendes Zeichen}
Was ist an einer Prüfung in Elektrotechnik fast immer ein gutes Zeichen?
{Auch WI-Studierende sitzen im Modul.}
{Das Oszilloskop startet sofort.}
{Die Aufgabenstellung enthält nur ein Bild.}
{Jemand sagt vor Beginn «das geht schnell».}
{A -- dann besteht zumindest Hoffnung, dass die Formeln nicht weglaufen.}

\mc{4.3 -- Die Digital-Engineering-Materialprobe}
Digital Engineering sagt: «Wir machen selbstverständlich auch Hardware.»
Welches Objekt wird als überzeugendster Beweis auf den Tisch gelegt?
{Ein Kabel, das tatsächlich an beiden Enden eingesteckt ist.}
{Ein Figma-Board mit sehr realistischem Schatten.}
{Ein QR-Code zur Dokumentation eines Kabels.}
{Ein Bluetooth-Symbol, aus Holz gelasert.}
{A -- ein funktionierendes Kabel ist im Zweifel das physischste Argument.}

\mc{4.4 -- Der gemeinsame Nenner}
Wann erkennen sich ET, WI und Digital Engineering ausnahmsweise sofort
gegenseitig als gleichwertige Menschen?
{Wenn das WLAN während einer Abgabe ausfällt.}
{Wenn jemand «Synergie» sagt.}
{Wenn ein Blockdiagramm drei Farben hat.}
{Wenn das Projekt einen Namen mit «Smart» trägt.}
{A -- Notlagen schaffen Interdisziplinarität.}

\mc{4.5 -- Projektstatus nach Studiengang}
Ein ET-Team sagt: «Es läuft.» Ein WI-Team sagt: «Es ist im Plan.» Ein
Digital-Engineering-Team sagt: «Die Experience ist stimmig.» Welche
Gegenfrage klärt alle drei Aussagen am schnellsten?
{«Kann man es messen?»}
{«Gibt es schon ein Logo?»}
{«Wer macht die Folien?»}
{«Ist die Farbpalette konsistent?»}
{A -- die Messung ist die diplomatische gemeinsame Sprache.}

\mc{4.6 -- Das gefährliche Kompliment}
Ein fachfremder Kollege sagt: «Das ist ja nur ein kleiner Sensor.» Was ist
die korrekteste ET-Antwort?
{«Klar. Bis du ihn anschliesst.»}
{«Dann hat er bestimmt keinen Offset.»}
{«Super, dann brauchen wir keine Masse.»}
{«Ein Sensor misst nur Gefühle.»}
{A}

\mc{4.7 -- Interdisziplinäre Definition}
Welche Aussage lässt ein ET-Team sofort vermuten, dass ein Projektmeeting
nun länger als das Projekt selbst dauert?
{«Lasst uns zuerst noch die Stakeholder abholen.»}
{«Wir haben die Stromaufnahme gemessen.»}
{«Der Kurzschluss ist lokalisiert.»}
{«Die Leiterplatte ist angekommen.»}
{A -- technisch nicht falsch, zeitlich aber selten harmlos.}

\mc{4.8 -- Die erste Frage an den Prototyp}
Ein Prototyp wird ausgestellt und leuchtet überraschend schön. Was fragt
ein ET-Mensch mit hoher Wahrscheinlichkeit als Erstes?
{«Was ist das für ein Strom?»}
{«Welcher Font ist das?»}
{«Wie ist die User Journey?»}
{«Kann man das skalieren?»}
{A -- Schönheit ist gut. Ein Strommesswert ist beruhigender.}

\runde{5. Der Teil, der in keiner Vorlesung vorkommt}

\mc{5.1 -- Das sauberste Laborprotokoll}
Welche Formulierung ist die ehrlichste?
{«Nach dem dritten Reset war das Verhalten reproduzierbar.»}
{«Der Aufbau funktionierte beim ersten Versuch fehlerfrei.»}
{«Die Masseklemme war nie ein Thema.»}
{«Das Datenblatt war überraschend eindeutig.»}
{A -- reproduzierbar ist hier ein mutiger, aber fachlich vertretbarer Begriff.}

\mc{5.2 -- Die natürliche Feindin des Prototyps}
Was bringt einen funktionierenden Prototypen am zuverlässigsten aus dem
Gleichgewicht?
{«Kannst du das rasch noch einmal vorzeigen?»}
{Eine zweite Messung}
{Eine beschriftete Leitung}
{Ein vollständig geladenes Netzteil}
{A -- das ist ein Feldversuch mit hoher Trefferquote.}

\mc{5.3 -- Das Labororakel}
Ein Messwert ist exakt das Gegenteil der Erwartung. Welche Erklärung ist
am wahrscheinlichsten?
{Der Massebezug ist irgendwo kreativ geworden.}
{Die Natur hat heute andere Konstanten.}
{Das Ohmsche Gesetz gilt nur bis Mittag.}
{Das Messgerät unterstützt jetzt Agile.}
{A -- danach darf man trotzdem noch das Messgerät prüfen.}

\mc{5.4 -- Der perfekte Kommentar}
Was ist die technisch korrekteste Reaktion auf den Satz «Ich habe nur
eine Kleinigkeit geändert»?
{«Welche Kleinigkeit, auf welchem Branch und mit welcher Versorgung?»}
{«Dann kann es nichts Wichtiges sein.»}
{«Perfekt, dann brauchen wir nicht zu messen.»}
{«Ich öffne schon einmal die Präsentation.»}
{A}

\mc{5.5 -- Ein System mit Persönlichkeit}
Wann ist eine Schaltung offiziell stabil?
{Wenn sie während der Demo nicht sichtbar schwingt.}
{Wenn die Polstellen gute Laune haben.}
{Wenn das Oszilloskop keine Fragen mehr stellt.}
{Wenn der Projektordner «final\_final» heisst.}
{A -- pragmatische Campusdefinition; die theoretische folgt im Zweifel später.}

\mc{5.6 -- Das letzte Wort im Labor}
Wer gewinnt eine Diskussion über ein elektrisches Verhalten am häufigsten?
{Die Messung}
{Die Person mit den meisten Folien}
{Die eleganteste Simulation}
{Derjenige, der «eigentlich» sagt}
{A -- manchmal sogar gegen die eigene Überzeugung.}

\offenfrage{5.7 -- Produktname}
Gib dem manuellen Gegenentwurf zur Autoscale-Taste einen Produktnamen.
{Juryentscheid. Gute Beispiele wären «Mettler-Mode», «Handskalierung Pro» oder
«Erst denken, dann drehen».}

\offenfrage{5.8 -- Abspann}
Vervollständige den Satz: «Der Prototyp war fast fertig, bis \ldots»
{Juryentscheid. Ein technischer Bezug ist Pflicht; ein plausibler Schuldiger
bringt Zusatzruhm.}

\vfill
\begin{center}
  \textit{Ende. Die richtige Antwort war selbstverständlich: erst messen.}
\end{center}

\end{document}
\fi
